\documentclass[11pt]{article}
\usepackage[margin=1in]{geometry}
\usepackage{amsmath,amssymb}
\usepackage{times}
\usepackage[colorlinks=true,linkcolor=blue,citecolor=blue,urlcolor=blue]{hyperref}
\usepackage{enumitem}
\usepackage{titlesec}

\titleformat{\section}{\normalfont\Large\bfseries}{\thesection}{1em}{}
\titleformat{\subsection}{\normalfont\large\bfseries}{\thesubsection}{1em}{}

\title{\textbf{Signal Before System: Whole-Body Imaging and the Construction of Organismal Activity}}
\author{Flyxion \\ \small Independent Researcher}
\date{September 2026}

\begin{document}

\maketitle

\begin{abstract}
Ruetten et al.\ (2026) report WHOLISTIC, a platform for imaging calcium dynamics across most major organ systems in the anterior body of a larval zebrafish at cellular resolution, paired with whole-body expansion microscopy and optogenetic perturbation. The result is presented, and has generally been received, as a step toward observing organismal physiology ``as a whole.'' This essay argues for a different framing. Whole-body observability does not disclose the organism directly. It constructs a sequence of admissible appearances, each produced by a distinct transform: pancellular sensor expression, optical acquisition, motion registration, segmentation, spectral clustering, anatomical attribution, and targeted intervention. Each stage preserves some distinctions and suppresses others. What the paper actually demonstrates is not the elimination of mediation but the disciplined binding of several partial witnesses into a claim that can survive independent tests. The essay develops this thesis through three extended cases: the registration step, which is normally treated as a technical preliminary but in fact performs a substantive classification of what counts as displacement versus signal; the paper's causal claims, which vary sharply in strength and should not be flattened into a single category of ``circuit discovery''; and the paper's bearing on distributed cognition, where a regulatory reading is defensible but a representational one is not licensed by the data. The paper is treated throughout as an unusually clean experimental object for sharpening, rather than merely illustrating, an existing set of distinctions concerning appearance, admissibility, and warranted inference.
\end{abstract}

\section{Introduction: the wrong way to praise this paper}

The natural response to Ruetten et al.\ (2026) is to describe it as a machine for seeing the whole organism at once. The paper's own language invites this: WHOLISTIC, an acronym for ``WHole-Organism Live-Imaging System for recording Tissue and IntraCellular activity,'' pairs pancellular expression of a calcium indicator with volumetric microscopy fast enough to track second-timescale dynamics across most of the anterior body, where the major viscera are concentrated, and follows the live imaging with whole-body expansion microscopy to recover molecular and ultrastructural context. Read quickly, the achievement sounds like the removal of a longstanding obstruction: physiologists have always known that organs talk to each other, and now, at last, someone has built the instrument that lets you watch the whole conversation.

That description is not false, but it flatters the method in a way that obscures what actually makes it valuable. No stage of this pipeline delivers the organism. Every stage delivers a rendering, produced under a specific transform, that is useful precisely because its limits are known and can be checked against other renderings produced under different transforms. The pancellular sensor decides in advance that calcium-dependent fluorescence, and only calcium-dependent fluorescence, will be the shared currency of ``activity.'' The optical system decides which depths, orientations, and volumes will be sampled and at what rate. The registration step decides which pixel-level changes count as an animal moving and which count as a cell changing state. Segmentation decides where one functional unit ends and another begins. Spectral clustering decides which temporal relationships count as organization. Anatomical attribution decides which cluster corresponds to which known or unknown cell type. Optogenetic perturbation decides, finally, whether a correlational pattern survives contact with a targeted intervention.

None of these decisions is illegitimate. All of them are necessary. But stacking them does not asymptotically approach a transparent view of ``the organism''; it produces a chain of admissible appearances, each of which can be interrogated, revised, or overturned by the others. The paper's real achievement is that the chain is unusually well instrumented: at almost every step, the authors provide the tools to ask whether the previous step's inference was warranted. This is the sense in which WHOLISTIC deserves to be called disciplined rather than holistic. It does not dissolve mediation. It manages it.

The remainder of this essay develops that claim through three cases drawn directly from the paper: the registration step, the paper's graded hierarchy of causal claims, and its implications for questions of distributed physiological organization and, at the far edge, distributed cognition.

\section{Registration as classification, not recovery}

Consider the transform most likely to be treated as a mere technical prerequisite: motion correction. A living, unrestrained set of organs inside a mounted but not fully immobilized larva will move. Muscle contractions, gut peristalsis, cardiac motion, and respiratory movement all displace tissue relative to the microscope's frame. If left uncorrected, a stationary calcium sensor moving in and out of a fixed voxel will produce an apparent fluorescence change that has nothing to do with calcium concentration. The paper addresses this with a registration system built on multiscale optical flow, aligning each acquired volume to a stable anatomical reference channel.

It is tempting to treat this step as bookkeeping: a necessary cleanup performed before the ``real'' analysis of calcium dynamics begins. That temptation should be resisted. Registration is itself an inferential act, and a consequential one. A more precise formulation makes this explicit. Let $M_t$ denote the organismal state at time $t$, and let the raw sensor output be
\[
x_t = R(M_t, \theta_{\mathrm{sens}}, \omega_t),
\]
where $\theta_{\mathrm{sens}}$ collects indicator kinetics, optical geometry, and detector properties, and $\omega_t$ collects noise, scattering, and uncorrected motion. The registered movie actually used in downstream analysis is not $x_t$ but
\[
\tilde{x}_t = \mathcal{A}_{\phi_t}\big(R(M_t, \theta_{\mathrm{sens}}, \omega_t)\big),
\]
where $\mathcal{A}_{\phi_t}$ is the alignment operator induced by an estimated deformation field $\phi_t$. The field $\phi_t$ is itself an inference, obtained under assumptions about local smoothness, patch size, choice of reference channel, and multiresolution search strategy. Nothing about those assumptions is dictated by the underlying physiology; they are chosen by the experimenters, and different reasonable choices would yield a different $\phi_t$.

Registration therefore imposes a practical partition on observed variation. Components explained by $\phi_t$ are treated as displacement and transformed away; components retained in $\tilde{x}_t$ remain available as candidate physiological signal; and components imperfectly represented by either model may be distributed across both or introduced through interpolation. The partition is neither binary at the level of the underlying organism nor lossless at the level of the recorded image. A fuller formulation makes the imperfection explicit:
\[
\tilde{x}_t
=
\mathcal{I}\!\left(
\mathcal{A}_{\widehat{\phi}_t}
\left[
R(M_t,\theta_{\mathrm{sens}},\omega_t)
\right]
\right)
+\varepsilon^{\mathrm{reg}}_t,
\]
where $\widehat{\phi}_t$ makes explicit that the deformation field is an estimate of bodily deformation rather than the deformation itself, $\mathcal{I}$ represents the interpolation and resampling that any practical alignment operation performs, and $\varepsilon^{\mathrm{reg}}_t$ collects the residual registration error left over once estimation and resampling are accounted for. This partition can fail in more than one direction at once. A deformation field fit too aggressively can absorb genuine physiological change, such as a contraction-linked shape change in a muscle fiber, and remove it from the signal before any biologist sees it. A deformation field fit too conservatively can leave true motion artifact inside $\tilde{x}_t$, where it will be indistinguishable from calcium-dependent fluorescence to every later stage of the pipeline, including spectral clustering, which has no way to know that a coherent pattern it has detected is actually a coherent pattern of incomplete motion correction rather than a coherent pattern of intracellular signaling. And some of what is lost or introduced will sit in neither $\phi_t$ nor $\tilde{x}_t$ cleanly, but in $\varepsilon^{\mathrm{reg}}_t$, where it is available to neither the motion model nor the physiological interpretation.

This matters for how the paper's discoveries should be read. Findings that involve tissue near sites of large, food- or peristalsis-driven deformation, such as the renal nephron waves or the multi-organ muscle synergies, inherit whatever residual uncertainty is present in $\widehat{\phi}_t$ for that tissue at that moment, in a way that findings involving relatively rigid or slow-moving structures, such as much of the hindbrain, do not. The paper does not ignore this. Its simultaneous anatomical reference channel, and its pre- and post-registration overlays and examinations of local alignment, directly test aspects of the inferred deformation field, while its Fourier temporal-correlation analysis separately characterizes which temporal frequencies the acquisition system can resolve at all. These address different failure modes and should not be treated as interchangeable validations: the former asks whether the estimated motion model is a good model of the motion, the latter asks what range of physiological dynamics the imaging rate and indicator kinetics can represent in the first place.

The nephron finding is the clearest case where this distinction earns its keep. It is among the results most plausibly affected by visceral displacement, given the concentration of gut and kidney motion in the imaged region, and it consequently inherits a registration burden greater than that of a signal observed in relatively rigid tissue. But its phase-ordered propagation and recurrence across functionally distinct compartments of the nephron are themselves positive evidence against a trivial motion explanation; a simple registration failure would not be expected to produce a structured, repeatedly propagating pattern with this specific organization. That evidence does not make the registration transform epistemically invisible, but it does mean the criticism should be read as inherited uncertainty rather than as grounds for suspicion. The point generalizes beyond this specific paper. Any claim that a pipeline has produced ``the corrected signal'' should be read as a claim that a particular, contestable partition of variation into displacement, signal, and residual has been performed and has, on available evidence, held up. It is an inscription, not a recovery, and its warrant is exactly as strong as the tests that have been run against it, no stronger.

\section{Grading the causal claims: association, convergence, necessity, sufficiency}

The paper reports several results that are often discussed under the single heading of ``circuit discovery.'' Collapsing them into one category understates how differently warranted they are, and risks letting the strength of the best case transfer, by association, to the weaker ones.

At one end of the paper's evidential spectrum sits the vagus--sphincter result. An unbiased whole-body coupling analysis first identifies a caudal hindbrain population whose activity tracks the gastropharyngeal sphincter. This is, on its own, an association: two signals covary, and covariation alone does not establish direction or mechanism. The paper then does something the whole-body screen by itself cannot do: it targets the implicated population with optogenetic stimulation and observes dose-dependent sphincter activation. This converts an association into evidence of sufficiency, at least for the specific manipulation performed. The paper goes further still, stimulating at successive anterior--posterior positions along the implicated population and observing correspondingly organized anterior--posterior effects on the viscera, which is evidence that the mapping is spatially systematic rather than coincidental. Finally, expansion microscopy shows that vagal innervation density falls off near the same anatomical boundary implicated functionally. Four largely independent lines of evidence, association, sufficiency under stimulation, spatial systematicity, and anatomical correspondence, converge on the same claim. Within this paper, this is the most complete causal case: association, intervention, spatial organization, and anatomical correspondence converge. It remains principally a sufficiency demonstration, however, rather than an exhaustive identification of the minimal and necessary pathway; a still stronger design would add pathway interruption, dose calibration, and temporal controls on top of the cell-specific recording and stimulation already performed. Its strength is also partly purchased by the tractability of the target. A sphincter controlled through a comparatively direct motor pathway is close to the limiting case in which a whole-body association can be reduced to a localized input--output relation. Success there demonstrates the platform's causal reach, but it should not be treated as representative of diffuse multi-organ coordination, and the paper is right to treat it as a flagship result precisely because it is the case where the platform's causal claims can be checked most exhaustively, not because every other finding in the paper is equally well supported.

The hypoxia and mesenteric blood flow result is weaker, though still substantial. Lowering ambient oxygen produces rapid constriction of the mesenteric artery and near-cessation of gut blood flow, while cerebral and muscular flow are comparatively preserved. Tricaine, a broad-spectrum anesthetic that suppresses neural activity and can independently affect cardiovascular physiology, abolishes the redistribution, showing that the response depends on a tricaine-sensitive physiological component; but the intervention is too broad to distinguish central neural control from peripheral neural, muscular, cardiac, or vascular effects, so it cannot by itself localize the necessary machinery to any specific circuit. Optogenetic inhibition of the hindbrain, which does restore mesenteric flow during continued hypoxia, provides the materially stronger localization and constitutes good evidence that intact hindbrain activity is necessary for the constriction under these experimental conditions. But necessity of the hindbrain for this outcome is not the same claim as identification of the hindbrain as the initiating oxygen sensor, nor does it establish that the implicated sympathetic neurons are the only output pathway, nor that gut hypoperfusion under hypoxia exists because it protects the brain, however plausible that adaptive story is. Each of those further claims would require its own targeted test.

Weaker still, though still informative, are the paper's coherence-based and lag-regression-based ``coupling'' findings, such as the cervical epaxial--abdominal muscle synergy and the general class of movement-associated responses across skin, gills, notochord, and ependymal cells. The authors are explicit that lag regression does not by itself establish causal direction, and this caveat should be taken seriously rather than treated as boilerplate. Common neural drive, mechanical coupling through the body wall, vasculature-mediated secondary effects, and imaging artifact from movement can all generate a signal that looks, to a clustering algorithm, exactly like a coupled response. The muscle synergy finding is strengthened by an independent anatomical argument, spino-occipital innervation capable of supporting joint recruitment of the two muscle groups, which is the right kind of corroboration to seek. But anatomical plausibility for a possible mechanism is not the same evidential category as an optogenetic perturbation that directly tests it, and the paper does not claim otherwise; readers summarizing the paper should not claim otherwise on its behalf.

The general point is that ``the platform can discover circuits'' is true only in the sense that the platform can generate candidate circuits and, in favorable cases such as the vagus--sphincter pathway, can carry a candidate through to a well-supported causal claim. It is not true that every coherence cluster the platform reports has been shown, or even claimed by the authors, to be a circuit in the causal sense. The section can be summarized as a rough evidential ordering,
\begin{align*}
\text{covariation}
\;<\;&\text{anatomical compatibility} \\
\;<\;&\text{interventional necessity or sufficiency} \\
\;<\;&\text{convergent pathway identification},
\end{align*}
though the ordering should not be read as a single absolute scale: anatomical compatibility and intervention answer different questions rather than merely different amounts of the same question, and a result can be strong on one axis while weak on another. What the ordering makes visible is why ``circuit discovery'' cannot be assigned at the first stage, covariation, on its own. A reader who wants to use this paper as evidence for a specific downstream claim should ask, for that specific claim, where on this ordering its support actually sits, rather than borrowing the confidence warranted by the paper's strongest case.

\section{Regulation without representation: the limit case for distributed cognition}

The paper's most provocative material, for anyone interested in questions of distributed organization and cognition, is the set of findings suggesting that non-neural tissue participates in temporally organized, body-wide coordination. Ependymal cells show calcium transients that track movement intensity during activity and shift into a slow, minutes-scale oscillatory regime during prolonged quiescence. The kidney nephron shows travelling, burst-like calcium waves rather than continuous activity. The vasculature reorganizes body-wide blood distribution under a single environmental perturbation, hypoxia, under direct hindbrain control. None of this activity is confined to, or even centered on, the nervous system in the conventional sense, even though several of the mechanisms the paper identifies are neurally gated.

It is worth being precise about what this licenses and what it does not. The findings are strong evidence for distributed regulation: tissues across the body sense local and systemic conditions, propagate that information through calcium signaling and, in several cases, through direct neural innervation, and adjust their own state and the state of neighboring tissue in response. This is already a meaningful correction to any account of organismal function that treats the nervous system as the sole locus of dynamically organized information processing and everything else as passive infrastructure executing its commands. The ependymal case is a particularly clean instance: expansion microscopy reveals long ventral processes reaching a region where a live calcium signal had already been detected but could not initially be reconciled with the conventionally described location of ependymal cell bodies. The apparent anatomical mismatch was not noise to be averaged away; it was a residual that forced revision of the structural picture, and the revision arrived from an entirely different modality than the one that produced the anomaly. That is a genuine demonstration that unexplained discrepancies between dynamic and structural data can be treated as information about a missing substrate rather than as error to be discarded.

None of this, however, establishes that peripheral tissues represent anything. Regulation is not representation. A tissue that senses local oxygen tension and constricts in response, or that propagates a calcium wave along its own length in a stereotyped rhythm, is doing something functionally important and dynamically organized, but it is not thereby encoding world-directed content, participating in a shared representational format with other tissues, or contributing a component to anything resembling a unified phenomenal field. Selective covariation with an environmental condition is not, by itself, sufficient to license a representational claim. At minimum, a representational interpretation needs some defensible account of what fixes the variable's content, some way of distinguishing misrepresentation from mere malfunction, and some account of how that content is subsequently used by another process rather than merely correlated with an outcome. WHOLISTIC measures selective responses and coordinated control; it does not test any of these further conditions, and nothing in the present results should be read as satisfying them by default. The paper does not claim otherwise, and its framing is generally careful on this point. The risk is downstream, in how such results get absorbed into broader arguments about extended or distributed cognition, where ``coordination'' and ``coupling'' can quietly slide into ``cognition'' and ``mind'' without any additional evidence being supplied for that specific, much stronger claim.

The paper's actual bearing on questions of a conscious or quasi-conscious distributed system, then, is narrower and more useful than the maximal reading suggests. What it undermines is the assumption that the nervous system is the only tissue in the body organized as a dynamically structured, feedback-coupled system operating on physiologically relevant timescales. What it does not do is supply evidence that this coupling constitutes, contributes representational content to, or is sufficient for anything like a unified experiential state. The more defensible bridge is that global physiological coupling of the kind demonstrated here is plausibly part of the substrate that constrains, maintains, and reconfigures the states available to a nervous system that is conscious for independent reasons, rather than being, itself, a distributed form of that consciousness. Under hypoxia, for instance, neural activity reorganizes circulation, altered circulation changes organ conditions, and those changed organ conditions feed back into the neural states that are subsequently available. The body is part of the state-transition geometry in which neural processes occur. That is a substantial claim, and this paper gives it real experimental content. It is a different, and much more modest, claim than the one sometimes implied by loose talk of body-wide coordination as itself a form of cognition.

\section{Conclusion: what the chain is for}

The paper's discovery-and-validation architecture is worth taking as a template independent of the specific biological findings. Broad, hypothesis-agnostic recording generates candidates. Cell-specific imaging narrows population identity. Expansion microscopy tests anatomical compatibility with the proposed mechanism. Optogenetic perturbation tests directional control. No single stage licenses a final mechanistic claim on its own, and the paper is at its best precisely where it declines to let one stage stand in for the whole chain, as in the vagus--sphincter case, and weaker where a single stage, such as a coherence cluster or a lag regression, is left to imply more than it can support on its own.

This is the sense in which the title of this essay should be read literally. Signal comes before system: the platform begins with measurements produced under specific, contestable transforms, and proceeds through a sequence of derived representations, registered movies, segmented components, functional clusters, anatomical attributions, whose biological interpretation must be earned rather than assumed at the moment they are produced. The organized system is not given in any one of these representations. It is the claim that becomes admissible when independent transformations preserve enough of the same structure to be checked against each other. The paper's genuine and considerable achievement is not that it finally lets anyone see the organism whole. It is that it assembles, with unusual explicitness, the chain of transforms required to move from raw, ambiguous, whole-body measurements to a claim that has been tested enough different ways to be believed, while leaving the seams between those transforms visible enough that a careful reader can see exactly how much weight each one is bearing.

\end{document}
